% Appendix: corpus documentation.
\section{Corpus and Evaluation Dataset}
\label{app:corpus}

% T1: the evaluated corpus by volume, with gold counts by cutoff.
\begin{table}[t]
\centering\small
\setlength{\tabcolsep}{2pt}
\begin{tabular}{@{}clrrrrr@{}}
\toprule
\nwhead{$t$} & \nwhead{Book} & \nwhead{Par.} & \nwhead{Words} & \nwhead{Ch.} & \nwhead{$|\Gt|$} & \nwhead{$|\Ft|$} \\
\midrule
1 & \emph{Anne of Green Gables} & \res{T1-anne-v1-par} & \res{T1-anne-v1-words} & \res{T1-anne-v1-ch} & \res{T1-anne-gold-t1} & \res{T1-anne-future-t1} \\
2 & \emph{Anne of Avonlea} & \res{T1-anne-v2-par} & \res{T1-anne-v2-words} & \res{T1-anne-v2-ch} & \res{T1-anne-gold-t2} & \res{T1-anne-future-t2} \\
3 & \emph{Anne of the Island} & \res{T1-anne-v3-par} & \res{T1-anne-v3-words} & \res{T1-anne-v3-ch} & \res{T1-anne-gold-t3} & \res{T1-anne-future-t3} \\
4 & \emph{Anne's House of Dreams} & \res{T1-anne-v4-par} & \res{T1-anne-v4-words} & \res{T1-anne-v4-ch} & \res{T1-anne-gold-t4} & \res{T1-anne-future-t4} \\
5 & \emph{Rainbow Valley} & \res{T1-anne-v5-par} & \res{T1-anne-v5-words} & \res{T1-anne-v5-ch} & \res{T1-anne-gold-t5} & \res{T1-anne-future-t5} \\
\midrule
 & Total & \best{\res{T1-anne-paras}} & \best{\res{T1-anne-words}} & \best{\res{T1-anne-chapters}} & & \\
\bottomrule
\end{tabular}
\caption{The evaluated corpus. Paragraphs, words and chapters are per volume; $|\Gt|$ is the
number of gold facts eligible at cutoff $t$ and $|\Ft|$ the number of later facts assessed for
disclosure at $t$.}
\label{tab:corpus}
\end{table}


\paragraph{Editions.} The Anne series is read from Project Gutenberg eBooks 45 (\emph{Anne of
Green Gables}), 47 (\emph{Anne of Avonlea}), 51 (\emph{Anne of the Island}), 544 (\emph{Anne's
House of Dreams}) and 5343 (\emph{Rainbow Valley}), in that order (Table~\ref{tab:corpus}).
The corpus follows the publication order of the first five Anne books (1908--1919), excluding
two later-written books set between them. Throughout the paper, a volume is one book; word
counts use whitespace-delimited tokens in the parsed text.

\paragraph{Development history.} The pipeline was developed on other series,
which contribute no reported results. The Anne books were processed while the pipeline was
still being revised, so this corpus is not a held-out evaluation. Baseline settings were chosen on George
Eliot's \emph{Middlemarch} (Gutenberg eBook 145, Books I--V) and not on the Anne series
(Appendix~\ref{app:baselines}).

\paragraph{Cast.} The evaluated cast contains the \res{T1-anne-cast} characters with at least one
gold fact. A comparison system generates a character page at cutoff $t$ if a gold fact about
that character is evidenced by $t$ or their full name occurs in $\Dt$. Characters without a
page remain in the evaluation denominators. This rule replaces an earlier rule that matched
any single word of a name and generated baseline pages for characters not yet introduced
(Appendix~\ref{app:casebook}).

\paragraph{Dataset size and units.}
The source corpus contains \res{T1-anne-words} whitespace-delimited words in
\res{T1-anne-paras} paragraphs and \res{T1-anne-chapters} chapters. The evaluated inventory
contains \res{T1-anne-gold-t5} selected facts about \res{T1-anne-cast} characters.
Table~\ref{tab:dataset-size} distinguishes these unique records from repeated judgments.
Each original or treated condition contains the same seven systems at five cutoffs;
renaming and paraphrasing reuse the inventory and reveal dates rather than adding an
independent corpus. Recall labels judge eligible fact--cutoff pairs; disclosure labels judge
later fact--cutoff pairs. Assertions are a separate system-generated population sampled
for precision and citation sufficiency. Repeated cutoffs share facts and characters.

\begin{table}[t]
\centering\small
\setlength{\tabcolsep}{4pt}
\begin{tabular}{@{}lrr@{}}
\toprule
\nwhead{Primary-matrix unit} & \nwhead{Original} & \nwhead{Treated} \\
\midrule
System--cutoff artifacts & \res{DATA-orig-artifact-cells} & \res{DATA-treated-artifact-cells} \\
Selected gold records & \res{T1-anne-gold-t5} & \res{T1-anne-gold-t5} \\
Cast characters & \res{T1-anne-cast} & \res{T1-anne-cast} \\
Recall judgments & \res{DATA-orig-recall-assessments} & \res{DATA-treated-recall-assessments} \\
Disclosure judgments & \res{DATA-orig-disclosure-assessments} & \res{DATA-treated-disclosure-assessments} \\
Sampled assertion records & \res{DATA-orig-precision-sample-rows} & \res{DATA-treated-precision-sample-rows} \\
Excluded frame errors & \res{DATA-orig-precision-frame-errors} & \res{DATA-treated-precision-frame-errors} \\
Evaluated precision atoms & \res{DATA-orig-precision-evaluated-atoms} & \res{DATA-treated-precision-evaluated-atoms} \\
\bottomrule
\end{tabular}
\caption{Size of the primary evaluation matrix, excluding earlier and prefix-build diagnostics.
Counts describe records and judgments, not independent trials. Both conditions reuse the same
gold inventory. Citation sufficiency uses the sampled assertion records; X1 supplies no citations.}
\label{tab:dataset-size}
\end{table}


\paragraph{What a record contains.}
The gold YAML groups facts by character. Each fact records a claim and a paragraph identifier;
the identifier encodes its evidence volume, chapter and paragraph. Optional keyword lists
assist disclosure search, while reader and note fields record provenance. The stored YAML
retains historical later-volume entries, but this evaluation selects evidence volumes
one to five only. Below is a minimal excerpt from an existing record, with provenance
metadata omitted:

\begin{lstlisting}[basicstyle=\ttfamily\scriptsize]
series: anne
characters:
  Matthew Cuthbert:
    facts:
      - claim: "Matthew is the shyest man alive."
        evidence: v01:c01:p0004
\end{lstlisting}

\paragraph{Easier example: a direct statement.}
The cited paragraph describes Matthew as ``the shyest man alive'' and says that he
``hated to have to go among strangers''. The human-accepted gold check
(Appendix~\ref{app:annotation}) labels this record fully supported. The character and description occur together, so matching
the proposition requires neither a cross-paragraph identity link nor a conjunction.
Its evidence volume is one: it is eligible at every evaluated cutoff. This is the
narrator's characterization of Matthew, rather than a claim about his current life status.

\paragraph{Harder example: conjunction and local coreference.}
Owen Ford's record reads: ``Owen Ford is a newspaper man and a grandson of the schoolmaster
who built the house'', with evidence \pid{v04:c22:p0019}. Miss Cornelia's report includes
``His name is Owen Ford'', ``a newspaper man'' and ``a grandson of the schoolmaster who
built this house''. The human-accepted gold check labels the record fully supported.
Interpreting the pronouns requires linking them to Owen; recall requires both the
occupation and the family/house relation. A generated statement that supplies only his
occupation does not convey the full record. At cutoff three the fact is a disclosure probe;
at cutoff four it becomes eligible. The example illustrates the existing conjunction
rule, not a measured difficulty grade or a newly assigned annotation.

\paragraph{An evidence boundary case.}
The gold claim ``Anne marries Gilbert Blythe'' points to \pid{v04:c04:p0000}, which begins
``Anne wakened on the morning of her wedding day''. That paragraph alone neither names
Gilbert nor establishes the completed marriage. The human-accepted gold check
therefore marks it partial, even though the fact is assigned to the correct volume.
This example distinguishes volume-level support from paragraph-level evidence sufficiency;
the recorded evidence defect remains part of the inventory, as reported in
Appendix~\ref{app:annotation}. These illustrative cases do not constitute new labels,
a measured easy/hard split, or a repair of the gold evidence.

\paragraph{Review population and release scope.}
The two-reviewer check covers \res{ANN-n-total} original-text items:
\res{ANN-n-gold} gold-evidence, \res{ANN-n-recall-sample} recall,
\res{ANN-n-precision-sample} precision and \res{ANN-n-m8-sample} citation checks, sampled
from the labels that existed when the sample was drawn (Appendix~\ref{app:annotation}).
Original and treated gold, per-system judgments, sampled assertions, candidate witnesses,
the two reviewers' final label files and recomputed agreement form the evaluation materials. The reviewer files, original response records and input
hashes are released; the human review is described in Appendix~\ref{app:annotation}.
